\section{Introduction}
For high-energy processes such as Higgs production at the Large-Hadron Collider,
quantum chromodynamics (QCD) factorization and parton distribution functions (PDFs) have been
essential for making theoretical predictions~\cite{Ellis:1991qj,Lin:2017snn}.
%
But for processes involving observation of a relatively small transverse momentum, $Q_\perp$
such as in Drell-Yan (DY) production and semi-inclusive deep inelastic scattering,
a new non-perturbative quantity called {\it soft function} is required  to capture
the physics of non-cancelling soft gluon-radiation at fixed $Q_\perp$~\cite{Collins:1981uk,Collins:1984kg,Ji:2004wu,Ji:2004xq}.
%
Physically, the soft function in DY is a cross section for a pair of a high-energy
quark and anti-quark (or gluon) traveling in the opposite light-cone directions
to radiate soft gluons of total transverse momentum $Q_\perp$ before they annihilate.
Although much progress has been made in calculating the soft function in perturbation
theory at $Q_\perp\gg \Lambda_{\rm QCD}$~\cite{Echevarria:2015byo,Li:2016ctv}, it is intrinsically non-perturbative when $Q_\perp$
is ${\cal O}(\Lambda_{\rm QCD})$.  Calculating the non-perturbative transverse-momentum-dependent (TMD)
soft function from first principles became feasible only recently~\cite{Ji:2019sxk}.


The main difference in such a calculation  in lattice QCD
is that it involves two light-like Wilson lines along directions
$n^\pm=\frac{1}{\sqrt 2}(1,\vec 0_\perp,\pm 1)$ in $(t,\perp,z)$ coordinates, making
direct simulations in Euclidean space impractical.  However,
much progress has been made in recent years in calculating physical quantities such as
light-cone PDFs using the framework of large-momentum effective theory (LaMET)~\cite{Ji:2013dva,Ji:2014gla}.
The key observation of LaMET is that the collinear  quark and gluon modes,
usually represented by light-like field correlators~\cite{Collins:2011zzd,Bauer:2000yr,Bauer:2001ct,Bauer:2001yt}, can be
accessed for large-momentum hadron states. A detailed review of LaMET
and its applications to collinear PDFs and other light-cone distributions
can be found in Refs.\cite{Ji:2020ect,Cichy:2018mum}. More recently,
some of the present authors have proposed that the TMD soft function
can be extracted from a special large-momentum-transfer form factor of
either a light meson or a pair of quark-antiquark color sources~\cite{Ji:2019sxk}.
Once calculated, the TMD factorization of the Drell-Yan and similar processes
can be made with entirely lattice-QCD-computable non-perturbative quantities~\cite{Ji:2014hxa,Ji:2018hvs,Ebert:2018gzl,Ebert:2019okf,Ji:2019ewn,Vladimirov:2020ofp}.


The TMD soft function is often defined and applied not in momentum space but in transverse coordinate space in terms of  the Fourier transformation variable $b_\perp $. In addition, it also depends on the ultraviolet (UV) renormalization scale $\mu$ (often defined
in dimensional regularization and minimal subtraction or $\overline{\rm MS}$) and rapidity regulators $Y+Y'$
\cite{Collins:2011zzd,Ji:2019sxk},
\begin{align}
S(b_\perp,\mu,Y+Y')=e^{(Y+Y')K(b_\perp,\mu)}S_I^{-1}(b_\perp, \mu)
\end{align}
where the first factor is related to rapidity evolution [described by the Collin-Soper (CS) kernel $K$], and the
second factor $S_I$ is the intrinsic, rapidity independent, part of the soft contribution.
The rapidity-regulator-independent CS-kernel $K$ is found calculable by taking ratio of
the quasi-TMDPDF at two different momenta~\cite{Ebert:2018gzl,Ebert:2019okf,Ji:2019ewn,Vladimirov:2020ofp,Ebert:2019tvc,Shanahan:2020zxr}.
On the other hand, calculating  the intrinsic soft function on the lattice
has never been attempted before.

In this paper we present the first lattice QCD calculation of the intrinsic soft function $S_I$
with several momenta on a 2+1 flavor CLS ensemble with $a=0.098$~fm~\cite{Bruno:2014jqa}, see Table~\ref{table:cls}. In particular we perform  simulations of the large-momentum light-meson form factor and quasi-TMD wave functions (TMDWFs), whose ratio gives the intrinsic soft function~\cite{Ji:2019sxk}. The Wilson loop matrix element will be  used  to remove the linear divergence in the quasi-TMD wave function. The CS kernel, $K$, can also be calculated from the external momentum dependence
of the quasi-TMD wave function~\cite{Ji:2020ect}, and
we will calculate it as a by-product. Our result is consistent with that of quenched lattice calculations of TMDPDFs~\cite{Shanahan:2020zxr}.
 
\begin{figure}[!th]
\begin{center}
\includegraphics[width=0.45\textwidth]{images/01_form_factor.pdf}
\caption{Illustration of the pseudo-scalar meson  form factor $F$ calculated in this work.
The initial and final momenta of the pion are large and opposite. The transition
``current'' is made of two local operators at a fixed spatial separation $b_\perp$. $t_{\rm sep}$
is the time separation  between the source and sink of the pion.}\label{fig:formfactor}
\end{center}
\end{figure}
